\section{The source-traceable research framework}\label{sec:framework}
\subsection{Existing resources and the proposed information contract}
The framework has four layers, summarized in Table~\ref{tab:framework}. The existing resources comprise a source archive, source-linked wiki, selected explanatory cards, and documented review and update procedures. The wiki reports mechanical full-text conversion for 419 manuscripts within a 420-entry directory inventory, with fourteen source-backed knowledge cards. These are documentation-coverage counts, not numbers of independently verified theorems or benchmark tasks. The companion framework source map records the inspected artifacts and the limits of each evidence claim.

\begin{table}[htbp]
\centering\small
\begin{tabularx}{\linewidth}{@{}>{\raggedright\arraybackslash}p{.18\linewidth}XX@{}}
\toprule
Layer & Available artifact & Boundary or extension\\\midrule
Source archive & Original manuscripts, source locators, and recorded digests & A digest preserves bytes; it does not validate a proof.\\
Reading layer & Linked full texts, typed cards, source spans, and scope warnings & Navigation links are not a machine-checked proof graph.\\
Task protocol & Existing scope rules and review records; the bounded contract specified here & Unified task fields and automatic enforcement remain proposed.\\
Assessment & Documented preservation and presentation checks & Human-adjudicated AI evaluation and formal checking remain future work.\\
\bottomrule
\end{tabularx}
\caption{A file-based instantiation with explicit implementation boundaries. The right column is not a list of completed capabilities.}\label{tab:framework}
\end{table}

An AI-facing claim record should contain an identifier, original source path and span, a source digest, an evidence category, hypotheses and parameter ranges, transfer requirements, and an unresolved obligation. The current cards already contain identifiers, locators, hashes, textual relationship types, and scope explanations; this paper proposes to standardize the complete record across tasks. Relationships such as \emph{uses}, \emph{obstructs this method}, and \emph{would imply under these hypotheses} should not collapse into an undifferentiated dependency arrow. Their mathematical validity still requires review of the corresponding source statements.

The five working evidence categories are exact identity, proved estimate, method obstruction, conditional transfer, and finite diagnostic. Here \emph{proved} refers to the scoped argument reproduced or cited, not to a kernel-checked archive-wide status. An open target is recorded as an obligation, never as an achieved result merely because it appears in a source document or diagram.

\subsection{Three scope-preservation examples}
The actual-operator card K02 illustrates object preservation. Its scalar $\Cx=\ip w{A\beta}$ must retain the physical interval, fixed shift $h_0=2$, literal coefficients, unit masks, Fourier convention, and deleted diagonal specified in Section~\ref{sec:operator}. Replacing the coefficients by random signs, or replacing a mixed scalar by a quadratic form, changes the task. Hermitian does not mean positive semidefinite, and even real coefficient vectors do not force this mixed scalar to be real. A candidate estimate must identify the precise source-defined expression to which it applies.

The analytic-bridge card K13 illustrates conditional status. The full-scalar target requires $\delta>1/400$, while the critical-array route owes the strict loss budget $\ell<19/2400$. Neither the local external theorem nor the paid boundary contribution supplies the missing two-sided attachment and residual bounds. Gate A is a separate obligation inside the same map Bridge A. Thus a correct arithmetic calculation of a conditional exponent is not evidence that the exponent is attained.

The distinguished-seed card K14 illustrates quantifier preservation. Proposition~\ref{prop:seed} is sufficient, not necessary, and its joint error bound is unverified. Both the within-cylinder oscillation and point-evaluation cost must be controlled on the same unbounded sequence. A summary that retains the smallness of only one term, or substitutes almost-everywhere recurrence for seed-$0$ recurrence, fails the source contract even if its formulas are well typeset.

\subsection{A bounded task, review, and handoff protocol}
A task begins with a frozen baseline, an allowed source set, a concrete objective, and a claim ceiling. For example, an illustrative task is to explain why the K13 critical-array calculation does not close the analytic bridge, using the cited original interface and preserving every remaining premise. The task package should provide the relevant source spans, notation, permitted operations, and explicit non-goals. Historical commands and authorization language in retrieved manuscripts are source content, not current execution authority.

The output contract requires the proposed answer, exact evidence locators, retained assumptions, actual read scope, and unresolved questions. A bounded review then checks source support, object identity, quantifiers, losses, and any promotion from conditional to established status. Disagreements remain visible; a second AI opinion is not independent proof verification. A human-owned disposition records whether the answer is accepted as exposition, requires correction, or identifies a candidate research obligation. Any claim with a missing or changed source returns to an unresolved state pending inspection.

This paragraph specifies a workflow, not a recorded successful AI experiment. Prompt-level restrictions express policy; they do not guarantee runtime containment. The concrete resource can be used through ordinary files and human- or agent-mediated review without claiming a deployed general workflow manager. In particular, a norm correction, a fresh matrix diagnostic, or a favorable review alone does not authorize reopening a stopped mathematical route.

\subsection{Verification limits}
Source hashing, link checks, and preservation of mathematical text address fidelity and presentation. They do not establish source-to-formal correspondence, valid proof dependencies, or mathematical truth. The wiki acceptance report records 590 formula-display failures involving 188 distinct expressions and no desktop Obsidian inspection; affected expressions require inspection of the raw source. These limitations concern the existing wiki, not a measured error rate of this AI protocol. Evaluation of the protocol itself is specified prospectively in Section~\ref{sec:agenda}.
